\documentclass[12pt]{article}
\usepackage{amsmath, amssymb}
\usepackage{geometry}
\usepackage{enumitem}
\usepackage{tcolorbox}
\usepackage{xcolor}

\geometry{margin=1in}

\tcbuselibrary{skins, breakable}
\newtcolorbox{answerbox}{
  colback=gray!10,
  colframe=gray!50,
  fonttitle=\bfseries,
  title=Answer,
  breakable
}

\title{\textbf{ESE 650: Learning in Robotics}\\
Problem 3 --- Practice Questions\\
(Quadtree, NeRF, Quaternions, SO(3), SE(3))}
\date{}

\begin{document}
\maketitle

% -------------------------------------------------------
\section*{Quadtree}

\subsection*{Q1}
In the domain $[0,1]\times[0,1]$ with finest cell size $[1/4, 1/4]$, insert the
following three points:
\[
  \{(0.1,\,0.1),\;(0.15,\,0.12),\;(0.9,\,0.8)\}.
\]
Draw the resulting quadtree. Compared with the exam problem (finest cell
$[1/8,1/8]$), why is the splitting depth different?

\begin{answerbox}
The first split divides $[0,1]^2$ into four $[1/2,1/2]$ cells.
$(0.1,0.1)$ and $(0.15,0.12)$ both fall in the lower-left quadrant; $(0.9,0.8)$
falls in the upper-right quadrant (one point only, so it never splits further).

The lower-left quadrant contains two points, so it splits again into four
$[1/4,1/4]$ cells.  At this resolution the two points land in different cells,
so splitting stops --- the finest cell size $[1/4,1/4]$ is already reached.

The depth is shallower than the exam problem because the coarser resolution is
sufficient to separate the two nearby points.
\end{answerbox}

% -------------------------------------------------------
\subsection*{Q2}
Using the original exam points
\[
  \{(0.52,\,0.75),\;(0.11,\,0.97),\;(0.1,\,0.98)\},
\]
change the finest cell size from $[1/8,1/8]$ to $[1/16,1/16]$.
Which nodes split further, and which do not?

\begin{answerbox}
$(0.52,0.75)$ is the only point in its cell at every resolution level, so it
\emph{never} triggers a split regardless of the finest cell size.

$(0.11,0.97)$ and $(0.1,0.98)$ are very close: $\Delta x = 0.01$,
$\Delta y = 0.01$.  At finest cell size $[1/8,1/8] = 0.125$, both points land
in the same minimum cell and splitting stops.  At $[1/16,1/16] = 0.0625$, they
still land in the same cell (since $0.01 < 0.0625$), so even with the finer
resolution the tree does \emph{not} split them into separate leaves.
\end{answerbox}

% -------------------------------------------------------
\section*{NeRF Position Encoding}

\subsection*{Q3}
The NeRF position encoding maps a coordinate $x$ to
\[
  \gamma(x)=\bigl[\cos(\omega_1 x),\,\sin(\omega_1 x),\,
                   \cos(\omega_2 x),\,\sin(\omega_2 x),\,\ldots\bigr].
\]
If position encoding is removed and the raw coordinates $(x,y,z)$ are fed
directly to the network, what problem arises? Explain using the concept of
\emph{spectral bias}.

\begin{answerbox}
Neural networks exhibit \textbf{spectral bias}: they preferentially learn
low-frequency functions and struggle to represent high-frequency ones.
A real scene contains sharp edges and fine textures --- high-frequency signals
--- that require the network to assign very different colors or opacities to
nearby 3-D locations.

Without position encoding the network receives only smooth, slowly varying
inputs, so it produces an overly blurred output regardless of its depth.
Position encoding explicitly embeds multiple frequency components into the input,
allowing the network to represent high-frequency scene content without requiring
an unreasonably large number of parameters.
\end{answerbox}

% -------------------------------------------------------
\subsection*{Q4}
NeRF takes both position $(x,y,z)$ and viewing direction $(\theta,\phi)$ as
inputs.

\medskip
\noindent\textbf{True or False:} Removing the viewing direction still allows
NeRF to correctly render specular (mirror-like) reflections.
Explain in 1--2 sentences.

\begin{answerbox}
\textbf{False.}  Specular reflection is \emph{view-dependent}: the same surface
point appears with different brightness and color depending on the observer's
direction (e.g.\ a highlight moves as you move).  Without the viewing direction
input the network can only assign a single fixed color to each point, making it
impossible to reproduce specular effects.
\end{answerbox}

% -------------------------------------------------------
\section*{Quaternions}

\subsection*{Q5}
The unit quaternion $q=(0,\,1,\,0,\,0)$ (i.e.\ $u_0=0,\,u_1=1,\,u_2=u_3=0$)
represents what 3-D rotation?  Describe it in axis--angle form.

\begin{answerbox}
Using the axis--angle parameterization
$q=\bigl(\cos\tfrac{\theta}{2},\;\sin\tfrac{\theta}{2}\,\hat{n}\bigr)$:
\[
  u_0 = \cos\tfrac{\theta}{2} = 0
  \;\Longrightarrow\;
  \tfrac{\theta}{2} = 90^\circ
  \;\Longrightarrow\;
  \theta = 180^\circ,
\]
and $\hat{n}=(1,0,0)$ (the $x$-axis).

Therefore $q=(0,1,0,0)$ represents a \textbf{rotation of $180^\circ$ about the
$x$-axis}.
\end{answerbox}

% -------------------------------------------------------
\subsection*{Q6}
\textbf{True or False:} Quaternion multiplication is commutative, i.e.\
$q_1\cdot q_2 = q_2\cdot q_1$.  Justify with an example.

\begin{answerbox}
\textbf{False.}  Quaternion multiplication encodes composition of 3-D rotations,
which is not commutative.  For example, rotating $90^\circ$ about the $x$-axis
\emph{then} $90^\circ$ about the $z$-axis yields a different final orientation
than performing the same rotations in the reverse order.
\end{answerbox}

% -------------------------------------------------------
\section*{SO(3)}

\subsection*{Q7}
A rotation matrix $R\in SO(3)$ has 9 scalar entries, yet $SO(3)$ is a
3-dimensional manifold.  Which constraints reduce the 9 degrees of freedom to 3?

\begin{answerbox}
$SO(3)$ is defined by two conditions:
\begin{enumerate}[noitemsep]
  \item $R^\top R = I$ (orthonormality)
  \item $\det R = +1$ (proper rotation, no reflection)
\end{enumerate}
The orthonormality condition $R^\top R=I$ imposes 6 independent scalar
constraints:
\begin{itemize}[noitemsep]
  \item Each of the 3 columns has unit norm: 3 constraints.
  \item Every pair of distinct columns is orthogonal: 3 constraints.
\end{itemize}
The determinant condition is then automatically satisfied (it rules out
$\det R=-1$) and does not add a further independent constraint.

Degrees of freedom: $9 - 6 = \mathbf{3}$.
\end{answerbox}

% -------------------------------------------------------
\section*{SE(3)}

\subsection*{Q8}
A robot starts at the origin.  It executes two steps (rotations and translations
measured in the \emph{body} frame):
\begin{align*}
  g_1 &= (R_1,\,t_1),\quad
        R_1 = \text{rotation of }90^\circ\text{ about }z\text{-axis},\quad
        t_1 = (1,0,0)^\top,\\
  g_2 &= (R_2,\,t_2),\quad R_2 = I,\quad t_2 = (1,0,0)^\top.
\end{align*}
Compute $g_2\circ g_1$ and find the robot's position in the world frame.

\begin{answerbox}
Using the composition rule $g_2\circ g_1=(R_2 R_1,\;R_2 t_1+t_2)$:

\textbf{Rotation:}
\[
  R_2 R_1 = I\cdot R_1 = R_1
  = \begin{pmatrix}0&-1&0\\1&0&0\\0&0&1\end{pmatrix}.
\]

\textbf{Translation:}
\[
  R_2 t_1 + t_2 = I\begin{pmatrix}1\\0\\0\end{pmatrix}
                 + \begin{pmatrix}1\\0\\0\end{pmatrix}
                 = \begin{pmatrix}2\\0\\0\end{pmatrix}.
\]

The robot is at world-frame position $\mathbf{(2,\,0,\,0)}$, facing the
direction given by $R_1$ (i.e.\ rotated $90^\circ$ about $z$).
\end{answerbox}

% -------------------------------------------------------
\subsection*{Q9}
\textbf{True or False:} Composition in SE(3) is commutative, i.e.\
$g_1\circ g_2 = g_2\circ g_1$.

\begin{answerbox}
\textbf{False.}  SE(3) inherits the non-commutativity of SO(3): because
rotations do not commute, the translation component $R_2 t_1$ in $g_2\circ g_1$
differs from $R_1 t_2$ in $g_1\circ g_2$ whenever $R_1\neq R_2$.
\end{answerbox}

% -------------------------------------------------------
\subsection*{Q10 (Synthesis)}
NeRF requires the camera pose of every training image.

\begin{enumerate}[label=(\alph*)]
  \item To which Lie group does a camera pose belong?
  \item What happens during NeRF training if some camera poses contain errors?
        Explain in 2--3 sentences.
\end{enumerate}

\begin{answerbox}
\textbf{(a)} A camera pose encodes both a 3-D rotation and a 3-D translation,
so it belongs to $\mathbf{SE(3)}$.

\textbf{(b)} NeRF fuses information from many images by projecting the same 3-D
point into each view.  If a pose is wrong, the projected pixel for that view
does not correspond to the correct 3-D location, so the network receives
contradictory color/density signals for the same point.  Unable to resolve the
contradiction, the network averages over the conflicting evidence and produces a
\emph{blurry} reconstruction.  This is why several methods jointly optimize
camera poses alongside the NeRF weights (\emph{pose refinement}).
\end{answerbox}

\end{document}
